% Abstract replacement paragraph for main-agent integration.
\begin{abstract}
PDF-based literature work can be slow because researchers need page-level evidence, citation support, and visible uncertainty before they can trust a draft answer. PaperMemory is a local-first compound AI assistant for postgraduate researchers who work with bounded PDF collections. It builds page-level EvidencePackets, combines visual and selectable-text retrieval, runs a bounded evidence-seeking loop, checks reliability, and keeps human verification in the workflow. Early market evidence from informal interviews with 10 PhD, master's, or postgraduate researchers gives a mean conditional purchase likelihood of 4.2/5 if the claimed functions are delivered, and 70\% prefer PaperMemory-provided API/model access through a monthly subscription rather than BYOK-only setup. We therefore evaluate a freemium model with a free BYOK tier and proposed 10 and 20~USD/month tiers that include capped model usage or API relay. On the local synthetic fixture suite, BM25 improves Recall@1 from 77.8\% to 88.9\% and retrieval-level refusal correctness from 33.3\% to 100.0\%; robustness fixtures cover missing evidence, citation drift, prompt-injection-like PDF text, conflicts, and low-text pages. A cost stress test estimates positive base-case net savings for first-pass evidence gathering under stated assumptions, while leaving conversion, retention, hosting cost, support cost, licensing, and real-corpus retrieval quality as open commercialization risks.
\end{abstract}

\section{Introduction}

Postgraduate researchers often spend early writing time turning PDFs into claims they can safely cite. The hard part goes beyond finding a relevant page. A useful research assistant must show why a page was used, which citation supports the sentence, and where the answer should stop because the evidence is thin. When that boundary is hidden, users still have to redo the evidence work by hand.

PaperMemory targets this paid workflow directly. In informal interviews with 10 university PhD, master's, or postgraduate researchers, the mean conditional purchase likelihood was 4.2/5 if PaperMemory delivered the promised functions: local PDF ingestion, evidence-backed answers, visible citations, bounded uncertainty, and a simpler model-access path. Seven of the 10 interviewees preferred PaperMemory-provided API/model access through a monthly subscription over a BYOK-only setup. This is early willingness-to-pay evidence, not product-market fit, but it is enough to justify testing a paid model in this course-stage report.

Existing research assistants such as Elicit and Consensus are priced around evidence-oriented workflows, although they are market comparators rather than executable baselines for this report \citep{elicitPricing,consensusPricing}. PaperMemory's narrower position is local-first PDF work: the user brings a scoped collection, the system builds page-level EvidencePackets, and the interface exposes citations, limits, and reliability status before the user turns the draft into research writing.

\begin{table}[t]
  \caption{A small postgraduate interview sample supports paid workflow interest, especially when PaperMemory provides managed model access, but the signal is early intent evidence.}
  \label{tab:market-validation}
  \centering
  \small
  \setlength{\tabcolsep}{4pt}
  \begin{tabularx}{\linewidth}{p{0.26\linewidth}p{0.18\linewidth}Y}
    \toprule
    Evidence item & Result & Boundary \\
    \midrule
    Interview sample & 10 postgraduate interviewees & Small interview sample, not a representative survey. \\
    Conditional purchase likelihood & 4.2 / 5 mean & Applies only if the claimed PDF, evidence, citation, and uncertainty workflow is delivered. \\
    API/model access preference & 70\% & Preference for product-provided API/model access through monthly subscription over BYOK-only setup. \\
    Proposed packaging & Free BYOK plus 10 and 20~USD/month tiers & Proposed revenue design, not actual conversion, retention, or production revenue. \\
    \bottomrule
  \end{tabularx}
\end{table}


The contribution is both technical and commercial. Technically, PaperMemory combines page rendering, selectable-text manifests, BM25 retrieval, page-canonical rank fusion, EvidencePacket validation, bounded evidence seeking, reliability checks, and an evidence-visible UI. Commercially, it connects that workflow to a freemium BYOK product and proposed paid tiers that remove API setup friction for users who prefer managed model access. The evaluation remains bounded: synthetic/local metrics and controlled live tests show implementation behavior, not broad scholarly-corpus superiority.

\section{Product Boundary and Revenue Logic}

The first target users are postgraduate researchers who work with dense PDF collections and need a faster first pass over evidence. A second group is small research, consulting, or R\&D teams that need traceable evidence packets before writing memos or reports. PaperMemory is not a systematic-review replacement or final expert adjudicator. It supports evidence gathering, citation packaging, and uncertainty surfacing before human review.

The proposed packaging follows the interview signal. The free tier is BYOK with limited groups, which lowers adoption friction and keeps model cost outside the product. The 10~USD/month tier adds PaperMemory-provided API/model access or API relay with capped included usage for routine postgraduate work. The 20~USD/month tier offers a larger workspace and higher capped included usage for heavier literature projects. The exact quotas should be set only after usage telemetry exists; this report intentionally does not invent them.

\begin{table}[t]
  \caption{The proposed freemium model connects interview preference for managed model access with bounded usage costs and evidence-workflow value.}
  \label{tab:revenue-model}
  \centering
  \small
  \setlength{\tabcolsep}{4pt}
  \begin{tabularx}{\linewidth}{p{0.18\linewidth}p{0.25\linewidth}p{0.25\linewidth}Y}
    \toprule
    Tier & User value & Model/API handling & Cost boundary \\
    \midrule
    Free BYOK & Limited groups for local PDF evidence work and adoption testing. & User supplies their own key or local model path. & No included model spend; product cost is mainly local workflow support. \\
    10~USD/month & Lower-friction postgraduate workflow with more groups and routine evidence runs. & Provider-compliant API relay for managed LLM access under a lower weekly cap, BYOK fallback, and rate limits. & Plausible only if managed usage stays bounded and high-use cases are throttled or upgraded. \\
    20~USD/month & Larger workspace for heavier literature projects and repeated report drafting. & Provider-compliant API relay under a higher weekly cap, larger workspace, and the same evidence controls. & Still requires fair-use policy, production hosting data, support-cost tracking, and licensing review. \\
    \bottomrule
  \end{tabularx}
\end{table}


The value logic is that users pay for convenience and evidence discipline together. The subscription removes API-key setup for the 70\% of interviewed users who preferred managed access, while the product keeps the local-first evidence workflow visible: pages, citations, limits, and reliability status remain part of the answer. That distinction matters because a plain hosted chatbot can be cheaper to build, but it does not give the user the same packet-level audit trail.

The cost model supports a plausible subscription experiment under bounded usage. The Node 9 stress test estimates API/compute costs of \$0.39 for a 20-PDF evidence packet, \$1.13 for a 50-PDF consulting/R\&D evidence scan, and \$1.62 for systematic-review pre-screening, with positive base-case net savings in all three scenarios after operation, verification, API/compute, and scenario license assumptions \citep{paperMemoryCostBenefit}. These numbers do not prove margins. They exclude production hosting, support, abuse controls, full licensing decisions, sales cycle cost, usage variance, actual conversion, and retention. The 10 and 20~USD/month tiers are therefore commercially plausible only under fair-use limits, and they remain proposed packaging rather than validated revenue.
